% Copy figures/memq_v9_paper_comparison_weighted.png to the report's images/
% directory before including this file after \appendix.

\section{Additional Analysis: Reproduction Fidelity and Zero-Shot Transfer}
\label{app:additional-analysis}

This appendix extends the reproduction experiments in Section~5 in two ways.
First, it separates the results reported in the original MemQ paper from the
results obtained by the reproduced v9 pipeline. Second, it evaluates whether
the fixed v9 model can transfer to GrailQA without any additional fine-tuning.
This distinction is important because a valid natural-language plan does not
necessarily imply that the memory module retrieves the correct Freebase
relations or reconstructs the correct executable query.

\subsection{Experimental Protocol}
\label{app:additional-protocol}

All additional experiments use the existing v9 planning model and the
direction fallback described in Section~5. No model parameters were updated.
WebQSP and CWQ are evaluated on the same test splits and local Freebase
snapshot as the main reproduction. For GrailQA, we use the labelled v1.0
development split \cite{guIIDThreeLevels2021}, because the official test
answers are not public. As in the original MemQ setting, the topic entity is
assumed to be known. The model input contains only the natural-language
question and the supplied topic-entity label. Gold relations, graph edges, and
answers are not included in the planning prompt.

The fixed v9 plan language cannot represent every GrailQA operator. We
therefore evaluate the largest directly compatible subset rather than silently
converting unsupported questions into ordinary path queries. Of 6,763
development examples, 4,970 (73.5\%) are retained. We exclude 1,251 examples
(18.5\%) that require functions such as counting, extrema, or comparisons, and
542 examples (8.0\%) containing literal constraints. Consequently, the
reported GrailQA result is a zero-shot transfer diagnostic on a supported
subset, not an official GrailQA leaderboard score. GrailQA++ is supported by
the same adapter, but no score is reported because an official labelled data
release was not available at the time of the experiment
\cite{duttGrailQAChallengingZeroShot2023}.

\subsection{Aggregation and Comparability}
\label{app:additional-aggregation}

The answer metrics use the same denominators as the executable-query
evaluation: 1,626 WebQSP questions, 3,273 CWQ questions, and all 4,970 retained
GrailQA questions. Graph metrics are computed for every successfully
reconstructed lookup item, yielding 1,635, 3,526, and 4,782 examples,
respectively. For a metric $m$ and hop groups $h$, the reported overall value
is weighted by the number of examples $n_h$ in each group:

\begin{equation}
    \overline{m}
    = \frac{\sum_h n_h\,\overline{m}_h}{\sum_h n_h}.
    \label{eq:weighted-hop-average}
\end{equation}

An unweighted mean over hop-level means would give a distorted result because
rare high-hop groups would receive the same influence as the much larger
one-hop group. For example, unweighted averaging of the seven populated
WebQSP Hits@1 hop bins gives 0.658, whereas the correct example-weighted value
is 0.822.

The graph-level comparison requires an additional qualification. The MemQ
paper reports EHR and GoldGED by hop in its Table~2, but it does not state
whether these values are computed on WebQSP, CWQ, or a combined sample
\cite{xuMemoryaugmentedQueryReconstruction2025}. Moreover, the original graph
matching implementation was not released. Our implementation normalizes
generated variable names before comparing relation-bearing edges and should
therefore be interpreted as a documented approximation rather than a bit-exact
reproduction of the paper's graph metrics. The paper values are reproduced in
Figure~\ref{fig:additional-paper-comparison} as reference rows, without
assigning them to either dataset.

\subsection{Aggregate Results}
\label{app:additional-results}

Table~\ref{tab:additional-aggregate-results} reports the aggregate results.
The original paper contains two relevant configurations. Its main comparison
in Table~1 reports 0.841/0.858 Hits@1/Macro-F1 on WebQSP and 0.803/0.830 on
CWQ. A separate analytical configuration in Table~3 reports slightly higher
values (0.857/0.872 and 0.817/0.845, respectively). These rows should not be
mixed when defining a single paper target. We use the main Table~1 values for
the direct comparison below.

\begin{table}[t]
    \centering
    \small
    \caption{Paper results and reproduced v9 results. EHR and GoldGED are
    reported only for the reproduction because the paper's hop-wise Table~2
    does not identify a dataset-specific aggregation.}
    \label{tab:additional-aggregate-results}
    \begin{tabular}{llrrrr}
        \toprule
        Dataset & Run & Hits@1 $\uparrow$ & Macro-F1 $\uparrow$
                & EHR $\uparrow$ & GoldGED $\downarrow$ \\
        \midrule
        WebQSP & Paper, Table 1 & 0.841 & 0.858 & -- & -- \\
               & Reproduced v9  & 0.822 & 0.823 & 0.502 & 1.783 \\
        \addlinespace
        CWQ    & Paper, Table 1 & 0.803 & 0.830 & -- & -- \\
               & Reproduced v9  & 0.723 & 0.738 & 0.497 & 2.648 \\
        \addlinespace
        GrailQA & v9 zero-shot, supported dev subset
               & 0.018 & 0.018 & 0.008 & 3.434 \\
        \bottomrule
    \end{tabular}
\end{table}

The WebQSP reproduction remains close to the paper on Hits@1 but is 3.5
percentage points lower in Macro-F1. The gap on CWQ is substantially larger:
8.0 points in Hits@1 and 9.2 points in Macro-F1. This agrees with the main
reproduction analysis: longer compositional queries, filters, entity binding,
and variable direction are more consequential on CWQ than on WebQSP. Hits@1
also depends on the first result returned by a \texttt{SELECT DISTINCT} query;
because SPARQL does not define an order without \texttt{ORDER BY}, this metric
can vary slightly across compatible Virtuoso deployments. Macro-F1 compares
answer sets and is therefore the more stable reproduction measure.

\begin{figure*}[t]
    \centering
    \includegraphics[width=\textwidth]
        {images/memq_v9_paper_comparison_weighted.png}
    \caption{Original MemQ paper values and the v9 reproduction by total hop
    count. Paper Table~2 provides hop-wise EHR and GoldGED but does not identify
    a dataset-specific aggregation. Paper Table~1 provides only dataset-level
    answer metrics. Reproduction averages are weighted over examples according
    to Equation~\ref{eq:weighted-hop-average}.}
    \label{fig:additional-paper-comparison}
\end{figure*}

\subsection{What the GrailQA Transfer Test Reveals}
\label{app:additional-grailqa-analysis}

The planning model generated a syntactically valid \texttt{Step1}-style plan
for all 4,970 retained GrailQA examples. Nevertheless, only 15 reconstructed
queries exactly matched the gold relation multiset, corresponding to a
structure accuracy of 0.003. EHR falls to 0.008, Hits@1 to 0.018, and Macro-F1
to 0.018. Performance also declines with depth: Macro-F1 is 0.020 for one-hop,
0.014 for two-hop, 0.005 for three-hop, and zero for the retained four-hop
examples.

This result separates two capabilities that can otherwise be conflated.
Fine-tuning teaches v9 to emit plans with the expected surface grammar, which
explains the 100\% plan-format validity. It does not provide a general mapping
from arbitrary natural-language relations to Freebase predicates. Query
reconstruction still depends on nearest-neighbour retrieval from a memory
constructed from WebQSP and CWQ relations. On GrailQA, a fluent plan can
therefore retrieve a semantically adjacent but graphically incorrect memory
entry. The very low one-hop EHR shows that the failure is not explained only by
longer reasoning chains; relation coverage and alignment of the plan language
with the memory embedding space already fail at the first edge.

The supported-subset design makes this conclusion stronger, not weaker. The
evaluation removes functions and literal constraints that v9 cannot express,
thereby favouring the model's existing path-query vocabulary. Despite this
advantage, transfer remains near zero. The result suggests that MemQ's strong
in-domain performance is based primarily on reuse of a closed relation memory,
not on dataset-independent query planning. Extending the method to a new
benchmark consequently requires at least rebuilding or expanding the query
memory, even if retraining the planning model is prohibited.

\subsection{Threats to Validity}
\label{app:additional-validity}

Four limitations constrain the interpretation of these results. First, the
local Freebase snapshot may differ from the snapshot used by the authors,
affecting query coverage and answer sets. Second, unordered SPARQL results make
Hits@1 partly endpoint-dependent. Third, the GrailQA number covers only the
73.5\% subset expressible by v9 and must not be presented as a full benchmark
score. Fourth, EHR and GoldGED are based on our documented graph-matching
approximation, while the paper's exact matching code and the dataset
aggregation behind Table~2 are unavailable. These limitations do not affect
the central qualitative observation: valid natural-language plans alone are
insufficient for zero-shot transfer when the query memory does not cover the
target relation distribution.
